\section{Thesis Structure}
\label{sec:introduction_thesis_structure}

\begin{figure}[t]
    \centering
    \resizebox{\textwidth}{!}{
        \begin{tikzpicture}[
            chapter/.style={rectangle, draw, rounded corners, minimum height=1.6cm,
                            minimum width=2.9cm, align=center, font=\small, text width=2.7cm},
            legend/.style={rectangle, draw, rounded corners, font=\small,
                           inner xsep=8pt, inner ysep=5pt},
            >={Stealth},
            % Each category has its own color AND its own border style, so the
            % figure stays readable in grayscale: dashed for technical
            % background, solid for surveys, double line for the method.
            background/.style={draw=myblue,  fill=myblue!12, dashed, line width=1pt},
            survey/.style={draw=mygreen, fill=mygreen!12},
            method/.style={draw=myred,   fill=myred!12, double=myred!12,
                           double distance=1.6pt},
        ]

        % Chapter number (small, on top) followed by the chapter name in bold.
        \newcommand{\chapterbox}[2]{{\footnotesize Chapter~\ref{#1}}\\[3pt]\textbf{#2}}

        % Two parallel tracks, deep learning (top) and reinforcement learning
        % (bottom), which converge on the method chapter.
        \node[chapter, background]  (dl)   at (0, 1.1)   {\chapterbox{chap:deep_learning}{Deep\\Learning}};
        \node[chapter, survey] (ls)   at (3.5, 1.1) {\chapterbox{chap:latent_space_analysis}{Latent\\Space}};
        \node[chapter, background]  (rl)   at (0, -1.1)  {\chapterbox{chap:reinforcement_learning}{Reinforcement\\Learning}};
        \node[chapter, survey] (xrl)  at (3.5, -1.1){\chapterbox{chap:explainability_in_reinforcement_learning}{Explainability\\in \gls{rl}}};
        \node[chapter, method]   (car)  at (7.5, 0)   {\chapterbox{chap:clustering_agent_representations}{\gls{car}\\Method}};
        \node[chapter, method]   (exp)  at (11, 0)    {\chapterbox{chap:experiments}{Experiments}};

        % Within each track.
        \draw[->, thick] (dl)  -- (ls);
        \draw[->, thick] (rl)  -- (xrl);

        % Both tracks converge on the method chapter.
        \draw[->, thick] (ls.east)  -- ++(0.55, 0) |- ([yshift=0.35cm]car.west);
        \draw[->, thick] (xrl.east) -- ++(0.55, 0) |- ([yshift=-0.35cm]car.west);

        \draw[->, thick] (car) -- (exp);

        \matrix[ampersand replacement=\&, column sep=0.8cm, inner sep=0pt] at (5.5, -2.8) {
            \node[legend, background]  {Technical background}; \&
            \node[legend, survey] {Survey}; \&
            \node[legend, method]   {Method}; \\
        };

        \end{tikzpicture}
    }
    \caption{Structure of the thesis.}
    \label{fig:thesis_map}
\end{figure}


The thesis is organized into a sequence of chapters, with each chapter falling into one of these three categories:
\begin{itemize}
\item[] \textcolor{myblue}{\textbf{Technical background}} chapters introduce fundamental concepts required to understand the subsequent contributions.
\item[] \textcolor{mygreen}{\textbf{Survey}} chapters review the existing literature around a specific research question.
The organization adopted in each of these surveys is original and therefore constitutes a contribution of this thesis.
\item[] \textcolor{myred}{\textbf{Method}} chapters present \gls{car}, our \gls{me} method for \gls{rl}.
\end{itemize}

Following these categories, the chapters are organized as shown in Figure~\ref{fig:thesis_map} and detailed below:
\begin{description}[style=nextline, leftmargin=1.5em, itemsep=0.5em]
\item[\textcolor{myblue}{\textbf{Deep Learning (Chapter~\ref{chap:deep_learning}, technical background)}}]
This chapter introduces the fundamental concepts of \gls{dl}, the training paradigm used to optimize deep \glspl{nn}.
It begins with an overview of the historical developments that led to the emergence of the field.
It then formalizes the objective of \gls{dl} and describes the structure of artificial \glspl{nn}.
Next, it explains how the learning process shapes these networks and how the resulting models are evaluated.
It concludes by showing that \glspl{nn} progressively transform their inputs into intermediate representations, whose study motivates the following chapter.

\item[\textcolor{mygreen}{\textbf{Latent Space Analysis (Chapter~\ref{chap:latent_space_analysis}, survey)}}]
This chapter addresses Research Question~\ref{rq:latent_space} by providing a survey of existing approaches for analyzing the \glspl{ls}.
It first makes explicit the working hypothesis that underlies much of this literature, which we refer to as the \ConceptHypothesisName.
It then reviews four families of methods, based respectively on probing, interpretable directions, feature superposition, and clustering.
These families are compared within a common synthesis, which highlights their shared goal of understanding how information is organized in \glspl{lr}.
The chapter closes by discussing the absence of consensus on how to evaluate whether these methods truly capture the concepts of a model.

\item[\textcolor{myblue}{\textbf{Reinforcement Learning (Chapter~\ref{chap:reinforcement_learning}, technical background)}}]
This chapter introduces the fundamental concepts of \gls{rl}.
It begins with an overview of the milestones that led to the emergence of \gls{rl} as a unified field.
It then clarifies what is meant by an agent and formalizes the problem that \gls{rl} algorithms aim to solve.
Next, it explains how a policy is progressively refined through interaction with the environment to approach the optimal policy.
It concludes by showing how the combination of \gls{rl} and \gls{dl} gives rise to agents that are both powerful and difficult to understand.

\item[\textcolor{mygreen}{\textbf{Explainability in Reinforcement Learning (Chapter~\ref{chap:explainability_in_reinforcement_learning}, survey)}}]
This chapter addresses Research Question~\ref{rq:taxonomy} by reviewing existing approaches to explainability in \gls{rl}.
It first clarifies what is meant by an explanation in the context of \gls{xai}.
It then presents a taxonomy that categorizes \gls{xrl} methods according to the subject of the explanation, the source of information and the form it takes.
This taxonomy is used to structure a state-of-the-art.
The chapter concludes with three observations from this review: the absence of explanations of the learning process, a lack of unified evaluation metrics, and few intrinsic methods for explaining the deep neural policies.

\item[\textcolor{myred}{\textbf{The CAR Method (Chapter~\ref{chap:clustering_agent_representations}, method)}}]
This chapter addresses Research Questions~\ref{rq:mechanisms} and~\ref{rq:evaluation} by proposing \gls{car}, a new method for explainability in \gls{rl} based on the analysis of \glspl{lr}.
It first lays out the theoretical assumption underlying the method, which we term the \RepresentationConvergenceHypothesisName.
It then details the method itself, which extracts \glspl{lr} from surrogate policies and automatically detects \glspl{cs} through clustering.
Next, it introduces the \gls{cu}, \gls{as}, and \gls{acu} metrics, which evaluate the consistency and the abstraction of the resulting clusters without manual intervention.
Finally, it situates \gls{car} within the broader landscape of \gls{me}.

\item[\textcolor{myred}{\textbf{Evaluation of the CAR Method (Chapter~\ref{chap:experiments}, method)}}]
This chapter presents the implementation and experimental evaluation of the proposed method.
It first describes the software implementation of the \gls{car} pipeline and the experimental configuration, based on six environments from the Atari benchmark.
It then evaluates how faithfully the surrogate policies reproduce the behavior of the initial policy.
Next, a quantitative analysis measures the consistency and the abstraction of the clusters across the layers of the network.
Finally, a qualitative analysis examines the semantic content of the clusters, which capture aspects such as episode progression, specific game situations, and behavioral patterns.

\item[\textbf{Conclusion (Chapter~\ref{chap:conclusion})}]
This chapter closes the thesis by revisiting each research question in light of the contributions presented.
It then discusses the limitations of these contributions.
These limitations open several directions for future research, which are presented next.
Finally, the chapter ends with a broader reflection on this work.
\end{description}

The mathematical notation used throughout this thesis is summarized on page~\pageref{chap:notation}, and the acronyms are listed on page~\pageref{chap:acronyms}.

Finally, in the interest of transparency, we specify how generative \gls{ai} tools were used in this work.
\Glspl{llm} served as assistants during the preparation of this thesis.
They were used for writing tasks, including translation, wording improvement, and proofreading.
They were also used for software development, including code implementation and review.
Finally, they were used to support scientific work, including the identification of relevant documents and discussion of approaches.
They were not used as autonomous scientific sources.
All generated outputs were systematically reviewed, and the author takes full responsibility for the content presented in this thesis.

Now that the research context and objectives have been presented, we can turn to the first chapter, which introduces the fundamentals of \gls{dl}.